\documentclass[10pt,a4paper]{amsart}
\usepackage{amsmath,amssymb,amsthm,graphicx}
\usepackage{verbatim}
\newenvironment{reference}{\comment}{\endcomment}
%\newenvironment{reference}{}{}
\newenvironment{slogan}{\comment}{\endcomment}
\newenvironment{history}{\comment}{\endcomment}

% For commutative diagrams we use Xy-pic
\usepackage[all]{xy}

% We use 2cell for 2-commutative diagrams.
\xyoption{2cell}
% Theorem environments.
%
\theoremstyle{plain}
\newtheorem{theorem}[subsection]{Theorem}
\newtheorem{proposition}[subsection]{Proposition}
\newtheorem{lemma}[subsection]{Lemma}

\theoremstyle{definition}
\newtheorem{definition}[subsection]{Definition}
\newtheorem{example}[subsection]{Example}
\newtheorem{exercise}[subsection]{Exercise}
\newtheorem{situation}[subsection]{Situation}

\theoremstyle{remark}
\newtheorem{remark}[subsection]{Remark}
\newtheorem{remarks}[subsection]{Remarks}

\numberwithin{equation}{subsection}

% Macros
%
\def\lim{\mathop{\mathrm{lim}}\nolimits}
\def\colim{\mathop{\mathrm{colim}}\nolimits}
\def\Spec{\mathop{\mathrm{Spec}}}
\def\Hom{\mathop{\mathrm{Hom}}\nolimits}
\def\Ext{\mathop{\mathrm{Ext}}\nolimits}
\def\SheafHom{\mathop{\mathcal{H}\!\mathit{om}}\nolimits}
\def\SheafExt{\mathop{\mathcal{E}\!\mathit{xt}}\nolimits}
\def\Sch{\mathit{Sch}}
\def\Mor{\mathop{\mathrm{Mor}}\nolimits}
\def\Ob{\mathop{\mathrm{Ob}}\nolimits}
\def\Sh{\mathop{\mathit{Sh}}\nolimits}
\def\NL{\mathop{N\!L}\nolimits}
\def\CH{\mathop{\mathrm{CH}}\nolimits}
\def\proetale{{pro\text{-}\acute{e}tale}}
\def\etale{{\acute{e}tale}}
\def\QCoh{\mathit{QCoh}}
\def\Ker{\mathop{\mathrm{Ker}}}
\def\Im{\mathop{\mathrm{Im}}}
\def\Coker{\mathop{\mathrm{Coker}}}
\def\Coim{\mathop{\mathrm{Coim}}}

% Boxtimes
%
\DeclareMathSymbol{\boxtimes}{\mathbin}{AMSa}{"02}

%
% Macros for moduli stacks/spaces
%
\def\QCohstack{\mathcal{QC}\!\mathit{oh}}
\def\Cohstack{\mathcal{C}\!\mathit{oh}}
\def\Spacesstack{\mathcal{S}\!\mathit{paces}}
\def\Quotfunctor{\mathrm{Quot}}
\def\Hilbfunctor{\mathrm{Hilb}}
\def\Curvesstack{\mathcal{C}\!\mathit{urves}}
\def\Polarizedstack{\mathcal{P}\!\mathit{olarized}}
\def\Complexesstack{\mathcal{C}\!\mathit{omplexes}}
% \Pic is the operator that assigns to X its picard group, usage \Pic(X)
% \Picardstack_{X/B} denotes the Picard stack of X over B
% \Picardfunctor_{X/B} denotes the Picard functor of X over B
\def\Pic{\mathop{\mathrm{Pic}}\nolimits}
\def\Picardstack{\mathcal{P}\!\mathit{ic}}
\def\Picardfunctor{\mathrm{Pic}}
\def\Deformationcategory{\mathcal{D}\!\mathit{ef}}
\usepackage{cleveref}
\crefname{originalreference}{original reference}{original references}
\Crefname{originalreference}{Original reference}{Original references}
\makeatletter
\newlabel{sites-modules-definition-morphism-locally-ringed-topoi}{{04HA}{source}{}{sourcepointer-1}{}}
\newlabel{sites-modules-definition-morphism-locally-ringed-topoi@cref}{{[originalreference][0][]04HA}{[1][0][]source}{}{}{}}
\makeatother
\crefname{originallemma}{original lemma}{original lemmas}
\Crefname{originallemma}{Original lemma}{Original lemmas}
\makeatletter
\newlabel{lemma-etale-site-locally-ringed}{{04KH}{source}{}{sourcepointer-2}{}}
\newlabel{lemma-etale-site-locally-ringed@cref}{{[originallemma][0][]04KH}{[1][0][]source}{}{}{}}
\makeatother
\crefname{originaltheorem}{original theorem}{original theorems}
\Crefname{originaltheorem}{Original theorem}{Original theorems}
\makeatletter
\newlabel{theorem-exactness-stalks}{{04K5}{source}{}{sourcepointer-3}{}}
\newlabel{theorem-exactness-stalks@cref}{{[originaltheorem][0][]04K5}{[1][0][]source}{}{}{}}
\makeatother
\crefname{originalreference}{original reference}{original references}
\Crefname{originalreference}{Original reference}{Original references}
\makeatletter
\newlabel{sites-modules-lemma-locally-ringed-morphism}{{04H9}{source}{}{sourcepointer-4}{}}
\newlabel{sites-modules-lemma-locally-ringed-morphism@cref}{{[originalreference][0][]04H9}{[1][0][]source}{}{}{}}
\makeatother
\crefname{originallemma}{original lemma}{original lemmas}
\Crefname{originallemma}{Original lemma}{Original lemmas}
\makeatletter
\newlabel{lemma-points-small-etale-site}{{04K6}{source}{}{sourcepointer-5}{}}
\newlabel{lemma-points-small-etale-site@cref}{{[originallemma][0][]04K6}{[1][0][]source}{}{}{}}
\makeatother
\crefname{originallemma}{original lemma}{original lemmas}
\Crefname{originallemma}{Original lemma}{Original lemmas}
\makeatletter
\newlabel{lemma-stalk-pullback}{{04K2}{source}{}{sourcepointer-6}{}}
\newlabel{lemma-stalk-pullback@cref}{{[originallemma][0][]04K2}{[1][0][]source}{}{}{}}
\makeatother
\crefname{originalreference}{original reference}{original references}
\Crefname{originalreference}{Original reference}{Original references}
\makeatletter
\newlabel{spaces-lemma-lift-morphism-presentations}{{02X1}{source}{}{sourcepointer-7}{}}
\newlabel{spaces-lemma-lift-morphism-presentations@cref}{{[originalreference][0][]02X1}{[1][0][]source}{}{}{}}
\makeatother
\crefname{originallemma}{original lemma}{original lemmas}
\Crefname{originallemma}{Original lemma}{Original lemmas}
\makeatletter
\newlabel{lemma-geometric-lift-to-cover}{{04JX}{source}{}{sourcepointer-8}{}}
\newlabel{lemma-geometric-lift-to-cover@cref}{{[originallemma][0][]04JX}{[1][0][]source}{}{}{}}
\makeatother
\crefname{originalreference}{original reference}{original references}
\Crefname{originalreference}{Original reference}{Original references}
\makeatletter
\newlabel{etale-cohomology-lemma-morphism-locally-ringed}{{04I5}{source}{}{sourcepointer-9}{}}
\newlabel{etale-cohomology-lemma-morphism-locally-ringed@cref}{{[originalreference][0][]04I5}{[1][0][]source}{}{}{}}
\makeatother
\crefname{originallemma}{original lemma}{original lemmas}
\Crefname{originallemma}{Original lemma}{Original lemmas}
\makeatletter
\newlabel{lemma-describe-etale-local-ring}{{04KF}{source}{}{sourcepointer-10}{}}
\newlabel{lemma-describe-etale-local-ring@cref}{{[originallemma][0][]04KF}{[1][0][]source}{}{}{}}
\makeatother
\title{Morphisms of algebraic spaces induce morphisms of locally ringed topoi}
\author{The Stacks Project Authors}
\date{stacks-project:spaces-properties.tex:local-snapshot-20261012; GNU FDL 1.2 or later}
\begin{document}
\maketitle
\noindent\textit{Editorial note.} All delivered mathematical text is the author's own, copied verbatim from the pinned Stacks Project chapter file spaces-properties.tex; no statement or proof step was rewritten, reordered or removed.\par\medskip
\noindent\textit{Editorial note.} Span recomputation by label pairing (LANE\_BRIEF 10.2): in this chapter the author's proof environment follows the statement environment as a standalone block; the spec covers the complete environment pair (statement span, proof span) as recorded in the pool candidate.\par\medskip
\noindent\textit{Editorial note.} Occupation scan: combined claim\_key + statement\_anchor + generic-line scan over all registration files (active and withdrawn) found 13 S10-geom published cards for spaces-properties.tex (prop-occupied.json, 2026-10-14); those 13 labels are excluded from this batch by design. For the 77 free labels: 0 occupied. The registry mentions field (7557 files) and focused-seeds (645) both point to the same S10-geom cards (dual-source check), no other occupation.\par\medskip
\ifdefined\hypertarget\hypertarget{sourcepointer-1}{}\fi
\noindent\textit{Source reference.} Original reference 04HA: sources/stacks-chow/tags.txt, source lines 5861--5861. Cross-chapter standard prerequisite cited by the author inside the counted span; it is used as a pointer and its own proof is not part of this item.\par\medskip
\ifdefined\hypertarget\hypertarget{sourcepointer-2}{}\fi
\noindent\textit{Source reference.} Original lemma 04KH: sources/stacks-chow/tags.txt, source lines 5973--5973. Cross-chapter standard prerequisite cited by the author inside the counted span; it is used as a pointer and its own proof is not part of this item.\par\medskip
\ifdefined\hypertarget\hypertarget{sourcepointer-3}{}\fi
\noindent\textit{Source reference.} Original theorem 04K5: sources/stacks-chow/tags.txt, source lines 5961--5961. Cross-chapter standard prerequisite cited by the author inside the counted span; it is used as a pointer and its own proof is not part of this item.\par\medskip
\ifdefined\hypertarget\hypertarget{sourcepointer-4}{}\fi
\noindent\textit{Source reference.} Original reference 04H9: sources/stacks-chow/tags.txt, source lines 5860--5860. Cross-chapter standard prerequisite cited by the author inside the counted span; it is used as a pointer and its own proof is not part of this item.\par\medskip
\ifdefined\hypertarget\hypertarget{sourcepointer-5}{}\fi
\noindent\textit{Source reference.} Original lemma 04K6: sources/stacks-chow/tags.txt, source lines 5962--5962. Cross-chapter standard prerequisite cited by the author inside the counted span; it is used as a pointer and its own proof is not part of this item.\par\medskip
\ifdefined\hypertarget\hypertarget{sourcepointer-6}{}\fi
\noindent\textit{Source reference.} Original lemma 04K2: sources/stacks-chow/tags.txt, source lines 5958--5958. Cross-chapter standard prerequisite cited by the author inside the counted span; it is used as a pointer and its own proof is not part of this item.\par\medskip
\ifdefined\hypertarget\hypertarget{sourcepointer-7}{}\fi
\noindent\textit{Source reference.} Original reference 02X1: sources/stacks-chow/tags.txt, source lines 3824--3824. Cross-chapter standard prerequisite cited by the author inside the counted span; it is used as a pointer and its own proof is not part of this item.\par\medskip
\ifdefined\hypertarget\hypertarget{sourcepointer-8}{}\fi
\noindent\textit{Source reference.} Original lemma 04JX: sources/stacks-chow/tags.txt, source lines 5953--5953. Cross-chapter standard prerequisite cited by the author inside the counted span; it is used as a pointer and its own proof is not part of this item.\par\medskip
\ifdefined\hypertarget\hypertarget{sourcepointer-9}{}\fi
\noindent\textit{Source reference.} Original reference 04I5: sources/stacks-chow/tags.txt, source lines 5891--5891. Cross-chapter standard prerequisite cited by the author inside the counted span; it is used as a pointer and its own proof is not part of this item.\par\medskip
\ifdefined\hypertarget\hypertarget{sourcepointer-10}{}\fi
\noindent\textit{Source reference.} Original lemma 04KF: sources/stacks-chow/tags.txt, source lines 5971--5971. Cross-chapter standard prerequisite cited by the author inside the counted span; it is used as a pointer and its own proof is not part of this item.\par\medskip
\begin{lemma}
\label{lemma-morphism-locally-ringed}
Let $S$ be a scheme.
Let $f : X \to Y$ be a morphism of algebraic spaces over $S$.
The morphism of ringed topoi $(f_{small}, f^\sharp)$
associated to $f$ is a morphism of locally ringed topoi, see
Modules on Sites,
Definition \ref{sites-modules-definition-morphism-locally-ringed-topoi}.
\end{lemma}

\begin{proof}
Note that the assertion makes sense since we have seen that
$(X_\etale, \mathcal{O}_{X_\etale})$ and
$(Y_\etale, \mathcal{O}_{Y_\etale})$
are locally ringed sites, see
Lemma \ref{lemma-etale-site-locally-ringed}.
Moreover, we know that $X_\etale$ has enough points, see
Theorem \ref{theorem-exactness-stalks}.
Hence it suffices to prove that $(f_{small}, f^\sharp)$
satisfies condition (3) of
Modules on Sites,
Lemma \ref{sites-modules-lemma-locally-ringed-morphism}.
To see this take a point $p$ of $X_\etale$. By
Lemma \ref{lemma-points-small-etale-site}
$p$ corresponds to a geometric point $\overline{x}$ of $X$.
By
Lemma \ref{lemma-stalk-pullback}
the point $q = f_{small} \circ p$ corresponds to the
geometric point $\overline{y} = f \circ \overline{x}$ of $Y$.
Hence the assertion we have to prove is that the induced map
of \'etale local rings
$$
\mathcal{O}_{Y, \overline{y}} \longrightarrow \mathcal{O}_{X, \overline{x}}
$$
is a local ring map. You can prove this directly, but instead we deduce it
from the corresponding result for schemes. To do this choose a commutative
diagram
$$
\xymatrix{
U \ar[d] \ar[r]_\psi & V \ar[d] \\
X \ar[r] & Y
}
$$
where $U$ and $V$ are schemes, and the vertical arrows are surjective
\'etale (see
Spaces, Lemma \ref{spaces-lemma-lift-morphism-presentations}).
Choose a lift $\overline{u} : \overline{x} \to U$ (possible by
Lemma \ref{lemma-geometric-lift-to-cover}).
Set $\overline{v} = \psi \circ \overline{u}$. We obtain a commutative
diagram of \'etale local rings
$$
\xymatrix{
\mathcal{O}_{U, \overline{u}} &
\mathcal{O}_{V, \overline{v}} \ar[l] \\
\mathcal{O}_{X, \overline{x}} \ar[u] &
\mathcal{O}_{Y, \overline{y}}. \ar[l] \ar[u]
}
$$
By
\'Etale Cohomology, Lemma \ref{etale-cohomology-lemma-morphism-locally-ringed}
the top horizontal arrow is a local ring map. Finally by
Lemma \ref{lemma-describe-etale-local-ring}
the vertical arrows are isomorphisms. Hence we win.
\end{proof}

\end{document}
